% GENERATED FILE. Edit canonical lessons, then run npm run generate:books.
% canonical-source-hash: fnv1a64:deaaff4e74b77080
% canonical-lessons: KA-C130-karadi, KA-C130-koti, KA-C130-mola, KA-C130-jinke, KA-C130-nari

\chapter{Bear, Monkey, Rabbit, Deer, Fox}
\label{ch:ka-bear-monkey-rabbit-deer-fox}
\hlchaptermodality{130}

\begin{chapteropening}
\textbf{By the end of this chapter:} \emph{I can say a bear, a monkey, a rabbit, a deer and a fox.}

Complete the last lesson of chapter 130: i can say a bear, a monkey, a rabbit, a deer and a fox.
\end{chapteropening}

\section[karaḍi]{\kn{ಕರಡಿ} (karaḍi) — a bear}

\label{lesson:KA-C130-karadi}



\begin{quote}
Before the new one: say the Kannada for a tiger, then the Kannada for a lion.
\end{quote}

\subsection*{You'll want to know: \kn{ಕರಡಿ}}
\textbf{\kn{ಕರಡಿ}} — \emph{karaḍi} — \textquotedblleft{}a bear\textquotedblright{}.

\begin{grammarlens}[title={What you've built}]
One more word for this chapter.
\end{grammarlens}

\subsection*{Your turn}
\begin{itemize}
\raggedright
  \item \emph{Say it:} \emph{karaḍi}
  \item \emph{Say it:} \emph{karaḍi}, once more
  \item \emph{Say it:} say \emph{karaḍi}
  \item \emph{Recall:} say the Kannada for a tiger, then the Kannada for a lion, then say \emph{karaḍi} again
\end{itemize}

\begin{tcolorbox}[breakable,colback=teal!4,colframe=teal!35!black,title={Before you move on}]
What does \kn{ಕರಡಿ} mean? (\textquotedblleft{}A bear\textquotedblright{}.) Say it once more.
\end{tcolorbox}
\section[kōti]{\kn{ಕೋತಿ} (kōti) — a monkey}

\label{lesson:KA-C130-koti}



\begin{quote}
Before the new one: say the Kannada for a lion, then the Kannada for a bear.
\end{quote}

\subsection*{You'll want to know: \kn{ಕೋತಿ}}
\textbf{\kn{ಕೋತಿ}} — \emph{kōti} — \textquotedblleft{}a monkey\textquotedblright{}.

\begin{grammarlens}[title={What you've built}]
One more word for this chapter.
\end{grammarlens}

\subsection*{Your turn}
\begin{itemize}
\raggedright
  \item \emph{Say it:} \emph{kōti}
  \item \emph{Say it:} \emph{kōti}, once more
  \item \emph{Say it:} say \emph{kōti}
  \item \emph{Recall:} say the Kannada for a lion, then the Kannada for a bear, then say \emph{kōti} again
\end{itemize}

\begin{tcolorbox}[breakable,colback=teal!4,colframe=teal!35!black,title={Before you move on}]
What does \kn{ಕೋತಿ} mean? (\textquotedblleft{}A monkey\textquotedblright{}.) Say it once more.
\end{tcolorbox}
\section[mola]{\kn{ಮೊಲ} (mola) — a rabbit}

\label{lesson:KA-C130-mola}



\begin{quote}
Before the new one: say the Kannada for a bear, then the Kannada for a monkey.
\end{quote}

\subsection*{You'll want to know: \kn{ಮೊಲ}}
\textbf{\kn{ಮೊಲ}} — \emph{mola} — \textquotedblleft{}a rabbit\textquotedblright{}.

\begin{grammarlens}[title={What you've built}]
One more word for this chapter.
\end{grammarlens}

\subsection*{Your turn}
\begin{itemize}
\raggedright
  \item \emph{Say it:} \emph{mola}
  \item \emph{Say it:} \emph{mola}, once more
  \item \emph{Say it:} say \emph{mola}
  \item \emph{Recall:} say the Kannada for a bear, then the Kannada for a monkey, then say \emph{mola} again
\end{itemize}

\begin{tcolorbox}[breakable,colback=teal!4,colframe=teal!35!black,title={Before you move on}]
What does \kn{ಮೊಲ} mean? (\textquotedblleft{}A rabbit\textquotedblright{}.) Say it once more.
\end{tcolorbox}
\section[jiṅke]{\kn{ಜಿಂಕೆ} (jiṅke) — a deer}

\label{lesson:KA-C130-jinke}



\begin{quote}
Before the new one: say the Kannada for a monkey, then the Kannada for a rabbit.
\end{quote}

\subsection*{You'll want to know: \kn{ಜಿಂಕೆ}}
\textbf{\kn{ಜಿಂಕೆ}} — \emph{jiṅke} — \textquotedblleft{}a deer\textquotedblright{}.

\begin{grammarlens}[title={What you've built}]
One more word for this chapter.
\end{grammarlens}

\subsection*{Your turn}
\begin{itemize}
\raggedright
  \item \emph{Say it:} \emph{jiṅke}
  \item \emph{Say it:} \emph{jiṅke}, once more
  \item \emph{Say it:} say \emph{jiṅke}
  \item \emph{Recall:} say the Kannada for a monkey, then the Kannada for a rabbit, then say \emph{jiṅke} again
\end{itemize}

\begin{tcolorbox}[breakable,colback=teal!4,colframe=teal!35!black,title={Before you move on}]
What does \kn{ಜಿಂಕೆ} mean? (\textquotedblleft{}A deer\textquotedblright{}.) Say it once more.
\end{tcolorbox}
\section[nari]{\kn{ನರಿ} (nari) — a fox}

\label{lesson:KA-C130-nari}



\begin{quote}
Before the new one: say the Kannada for a rabbit, then the Kannada for a deer.
\end{quote}

\subsection*{You'll want to know: \kn{ನರಿ}}
\textbf{\kn{ನರಿ}} — \emph{nari} — \textquotedblleft{}a fox\textquotedblright{}.

\begin{grammarlens}[title={What you've built}]
One more word for this chapter.
\end{grammarlens}

\subsection*{Your turn}
\begin{itemize}
\raggedright
  \item \emph{Say it:} \emph{nari}
  \item \emph{Say it:} \emph{nari}, once more
  \item \emph{Say it:} say \emph{nari}
  \item \emph{Recall:} say the Kannada for a rabbit, then the Kannada for a deer, then say \emph{nari} again
\end{itemize}

\begin{tcolorbox}[breakable,colback=teal!4,colframe=teal!35!black,title={Before you move on}]
What does \kn{ನರಿ} mean? (\textquotedblleft{}A fox\textquotedblright{}.) Say it once more.
\end{tcolorbox}
